\documentclass[12pt]{article}
\usepackage[margin=1in]{geometry}
\usepackage{amsmath}
\usepackage{amssymb}
\usepackage{amsthm}
\usepackage{enumitem}
\usepackage{xcolor}
\usepackage{pagecolor}
\usepackage{fancyhdr}
\usepackage{bm}
\usepackage{etoolbox}
\usepackage{mdframed}

\definecolor{cream}{HTML}{faf7f1}
\pagecolor{cream}

% Contradiction symbol
\newcommand{\contradiction}{\(\lightning\)}
\usepackage{stmaryrd}
\newcommand{\contra}{\ensuremath{\Rightarrow\Leftarrow}}

\newtheoremstyle{proofstyle}
  {}{}{}{}{\bfseries}{.}{ }{}
\theoremstyle{proofstyle}

% Case indentation environment
\newenvironment{caseblock}[1]{%
  \begin{quote}
  \noindent\textbf{#1}\ %
}{%
  \end{quote}
}

\makeatletter
\renewcommand{\maketitle}{%
  \begin{flushleft}
    {\Large\bfseries CS173: Discrete Structures}\\[0.3em]
    {\Large\bfseries Proving induction}\\[0.5em]
    {\normalsize Mohammed Al Salhi}
  \end{flushleft}
  \vspace{1em}
}
\makeatother

\AtBeginDocument{\mathversion{bold}}

\begin{document}
\maketitle

\noindent\textbf{Theorem (Weak Induction).} For any wff $\varphi$ with one free variable, the following equivalence holds:
\[
(\forall n \in \mathbb{N})(\varphi(n)) \iff \varphi(0) \land (\forall k \in \mathbb{N})(\varphi(k) \Rightarrow \varphi(S(k)))
\]

\noindent\textbf{Theorem.}
\[
\forall M\bigl((M \subseteq \mathbb{N} \land M \neq \emptyset)
\Rightarrow \exists m\bigl(m \in M \land (\forall n \in M)(m \leq n)\bigr)\bigr)
\]

\begin{proof}[Proof]

\noindent$(\Leftarrow)$ Assume $\varphi(0)$ and $(\forall k \in \mathbb{N})(\varphi(k) \Rightarrow \varphi(S(k)))$.

Let $n \in \mathbb{N}$. Towards a contradiction, assume $\neg\varphi(n)$.

Consider $M := \{m \mid m \in \mathbb{N} \land \neg\varphi(m)\}$. Observe $M \subseteq \mathbb{N}$.

Observe $n \in M$ (because $n \in \mathbb{N}$ and $\neg\varphi(n)$). Therefore $M \neq \emptyset$.

So, there exists $m \in M$ such that $(\forall x \in M)(m \leq x)$.

\begin{caseblock}{Case 1: Suppose $m = 0$.}
  Recall $\varphi(0)$, so $\varphi(m)$. However, $\neg\varphi(m)$ because $m \in M$. \contra
\end{caseblock}

\begin{caseblock}{Case 2: Suppose there exists $w \in \mathbb{N}$ such that $S(w) = m$.}
  This means $w \cup \{w\} = m$. So $w \in m$. So $w < m$.

  \begin{caseblock}{Case 2.1: Suppose $\neg\varphi(w)$.}
    Then $w \in M$. Since $(\forall a \in M)(m \leq a)$, we have $m \leq w$.\\
    By definition, $m \in w$ or $m = w$. However, $m \neq w$.\\
    So $m \in w$. However, $\forall x\, \forall y\,(x \in y \Rightarrow y \notin x)$. \contra
  \end{caseblock}

  \begin{caseblock}{Case 2.2: Suppose $\varphi(w)$.}
    Then, by assumption, we know $\varphi(w) \Rightarrow \varphi(S(w))$.\\
    So $\varphi(S(w))$. So $\varphi(m)$. However, $\neg\varphi(m)$. \contra
  \end{caseblock}
\end{caseblock}

\medskip
\noindent$(\Rightarrow)$ Assume $(\forall n \in \mathbb{N})(\varphi(n))$.

Since $0 \in \mathbb{N}$, we know $\varphi(0)$.

Let $k \in \mathbb{N}$. Assume $\varphi(k)$.

Since $k \in \mathbb{N}$, we know $S(k) \in \mathbb{N}$. Therefore $\varphi(S(k))$.

\end{proof}

\bigskip
\noindent\textbf{Theorem.}
\[
(\forall n \in \mathbb{N})(1 + n = S(n))
\]

\begin{proof}[Proof]
By weak induction on $n$.

\medskip
\noindent\textbf{Base Case:} Observe
\[
1 + 0 = 1 \quad \text{by definition of addition}
\]
\[
= S(0) \quad \text{by definition of } 1
\]

\medskip
\noindent\textbf{Inductive Step:} Let $k \in \mathbb{N}$.

\noindent We need to show (NTS): $1 + S(k) = S(S(k))$, assuming the inductive hypothesis $1 + k = S(k)$.

\medskip
Assume $1 + k = S(k)$. Observing:
\[
1 + S(k) = S(1 + k) \quad \text{by definition of addition}
\]
\[
= S(S(k)) \quad \text{by the inductive hypothesis}
\]

Therefore, $(\forall n \in \mathbb{N})(1 + n = S(n))$.

\end{proof}
\end{document}